%% kapa: exam=1 round=33 year=2022 period=1 form=A subject=economics subject_ko=경제학원론 seq=16 number=56
%% source: Q-Net 감정평가사 기출문제 (cst00309) artlSeq=5209997 file=raw/1008/5209997/2022년도 제33회 감정평가사 1차시험 1교시(민법,경제학,부동산학).pdf
%% conversion: kordoc 4.24.51 -> tools/gen1.py (auto); review=proofread-sonnet; flags=none
\begin{question}{56}
두 재화 $X$와 $Y$만을 소비하는 두 명의 소비자 갑과 을이 존재하는 순수교환경제에서 갑의 효용함수는 \\
$U_{\text{갑}}(X_{\text{갑}}, Y_{\text{갑}}) = \min\{X_{\text{갑}}, Y_{\text{갑}}\}$, 을의 효용함수는 $U_{\text{을}}(X_{\text{을}}, Y_{\text{을}}) = X_{\text{을}} \times Y_{\text{을}}$이다. 갑과 을의 초기 부존자원($X$, $Y$)이 각각 (30, 60), (60, 30)이고 $X$재의 가격이 1이다. 일반균형(general equilibrium)에서 $Y$재의 가격은?
\begin{choices}
  \item 1/3
  \item 1/2
  \item 1
  \item 2
  \item 3
\end{choices}
\end{question}
\answer{3}
